\section{Problem Formulation}
\label{sec:problem_formulation}
% \subsection{System Model}

In this paper, we consider a set $\mathcal{N}$ of $N$ autonomous agents moving within a given area of space. We consider a set of $R$ \textit{regions of interest} in a given area, where $R_k$ is the $k$-th element of the set. We assume each region is a sphere with an equal radius, however, the formulation could extend to more complex models.
Without loss of generality, we assume agents move along an undirected graph $G = (V, E)$ that is placed in the same space as the regions of interest. Each vertex $v_i \in V$ represents a location, and each edge $e_{ij}$ represents a feasible route from vertex $v_i$ to $v_j$. A key property of this graph is that it traverses at least once every region of interest.
This graph typically models constraints to agent trajectories due to the morphology of the monitored environment, presence of obstacles, regions of interest distribution, and characteristics of agent movements, among others. The specific way in which the graph is derived is thus application and context-dependent  \citep{nguyen2022multi}. The graph is defined at the beginning of the mission, it does not change over time and it is known by all agents.

The \emph{path} of agent $n$, denoted as $p^n$, is an ordered list of edges $p^n = (e^n_1, e^n_2,...)$, such that two adjacent edges in the path are connected by a vertex of the graph. With $p = (p^1, ..., p^n, ..., p^N)$ we denote the joint paths of every agent. Let $B$ denote the maximum path length of each agent, which equals the number of edges an agent can traverse. Such a maximum value is derived from the agent's speed, but it may also capture various constraints, e.g. due to finite storage capacity, among others. To any path $p^n$ we associate a cost $b(p^n)$, equal to the number of traversed edges. A region $R_k$ is \textit{observed} if it is traversed by a path of an agent. Every region $R_k$ is associated with a utility $U(R_k)$, which models the value of the information that agents may collect from it. For ease of analytical treatment, we assume that it takes one unit of time to traverse any edge and that any exchange of information among agents is instantaneous. Note however that our approach can be easily extended to account for nonzero exchange duration, as well as for edge traversal times that differ among edges and agents.

Finally, we assume each agent can exchange information at any point in time with any other agent. This models scenarios in which agents have a wireless interface to a cellular access network. We assume the information exchange to be instantaneous, independent from the amount of information shared, and reliable, with no loss. In the experimental section, we relax this assumption and investigate the impact of nonidealities in information sharing on the effectiveness of our approach.

\subsection{Multi-agent planning with attrition}

\label{sec:attrition_settings}

We denote the information-gathering task as a \textit{mission}, for which each agent performs independent actions to achieve a collective goal - maximizing the global utility for the whole team. Each agent $n$ plans its path $p^n$ and coordinates with others in a decentralized manner while satisfying the given budget constraint $B$ on path length.
This formulation of the information gathering problem generalizes many multi-agent path planning problems, such as team orienteering problem~\citep{best2018online}.
%\footnote{\GR{Please add again here the formal definition of the information gathering problem you once put here.}}
We consider a scenario, in which mission planning is performed in a decentralized manner for scalability and computational feasibility, as mentioned in Section 1.
%\HN{Need to motivate a decentralized approach first}\footnote{I see better such a motivation in the intro section, not here.}, 
Thus, each agent plans its path while considering the potential actions of other agents and the team's total utility.
%\HN{I thought we have decided to completely remove communication impact and always assume full comms? I would remove this whole paragraph}
%We are now ready to describe the multi-agent in attrition settings. 
We consider scenarios where a subset $\mathcal{F}$ of the N agents fail during the mission. We focus on \textit{hard} failures, where agents interrupt reward collection and information exchange.
We assume the set of agents that fail $\mathcal{F}$ is unknown in advance and the time at which they fail to be determined by any arbitrary criteria or distribution. Therefore, our solution does not rely on knowing its size and probability distribution.
In the occurrence of a failure, all the utility collected by the failed agent is lost, i.e. it is not considered anymore in the computation of the global utility of the mission. This models a typical setup in information gathering, in which data collected by agents is relayed to data sinks only at the end of the mission.

Our goal is to provide an efficient planning and coordination mechanism that can quickly adapt to agent failures and maximize the global utility of the mission within the agent's budget constraint. Let $\mathcal{P} = (P^1, ..., P^n, ..., P^N)$, with $P^n$ denote the set of all possible paths of length $B$ which starts at agent $n$ starting position.
We define the following problem:

\begin{problem}\label{prob:1}(\textit{Multi-agent planning in attrition settings)}
\begin{equation}\label{eq:objectivefunction}
% \underset{p \in \mathcal{P}}{\text{maximize}}\   U_g(p) = \sum_{n \in \mathcal{N} \setminus \mathcal{F}} U(p^n)
\underset{p \in \mathcal{P}}{\text{maximize}}\   U_g(p)
\end{equation}
	\begin{align}
           \text{Subject to:} &  \quad b(p^n)\leq B,  \quad \forall n \in \mathcal{N}   \label{constraint_pathlength} \\
           & \quad 0 \leq | \mathcal{F} | \leq N
	\end{align}
\end{problem}

Constraint (\ref{constraint_pathlength}) derives from imposing that the total path length for the agent $n$ to be less than the travel budget $B$ available to each agent. Intuitively, our goal is to find a path for each agent such that the global utility associated with all regions observed by all agents during the mission is maximized, while a subset $\mathcal{F}$ of agents fail. Such an optimization problem cannot be solved efficiently. Indeed it is easy to see that Problem 1 is a variant of the well-known NP-hard traveling salesman problem.
%In the next section, we provide an adaptive and distributed planning solution to Problem 1, which is based on Monte Carlo Tree Search in combination with Regret Matching, that can quickly and efficiently adapt the plan in attrition settings.
% \GR{Introduce what follows: in the following section, we propose.... (relate what follows to the problem formulation: do we solve the problem? or what?}

% In the next section, we describe the state-of-the-art approach for the multi-agent active information gathering problem, namely Decentralized MCTS~\citep{best2019dec}, and explain why it fails in solving Problem 1.







